% Build with: tectonic paper.tex
\documentclass[11pt,letterpaper]{article}
\usepackage[T1]{fontenc}
\usepackage{lmodern}
\usepackage[margin=1in]{geometry}
\usepackage{amsmath,amssymb,amsthm,mathtools}
\usepackage{microtype,booktabs,hyperref}
\hypersetup{hidelinks,
 pdftitle={Improved Bounds for Depth-Two Threshold Circuits Computing Andreev's Function},
 pdfauthor={Christian Reitwiessner}}
\newtheorem{theorem}{Theorem}
\newtheorem{lemma}[theorem]{Lemma}
\newtheorem{corollary}[theorem]{Corollary}
\newtheorem{proposition}[theorem]{Proposition}
\theoremstyle{remark}
\newtheorem*{remark}{Remark}
\newcommand{\bits}{\{0,1\}}
\newcommand{\ind}[1]{\bigl[\,#1\,\bigr]}
\newcommand{\THR}{\mathsf{THR}}
\newcommand{\cH}{\mathcal H}
\title{Improved Bounds for Depth-Two Threshold Circuits Computing Andreev's Function}
\author{Christian Reitwiessner}
\date{November 2026}
\begin{document}
\maketitle

\begin{abstract}
Andreev's function on $n$ input bits reads a bit from a
table at an address given by the parities of blocks of input bits. Kane and
Williams used it for the first superlinear gate lower bounds against
depth-two threshold circuits and asked for tight bounds. We improve the bounds by showing:
\begin{itemize}
\item an exact circuit with $(5/16+o(1))n^2$ gates and polynomial integer
      weights (previously $O(n^3/\log n)$);
\item for agreement on a $\tfrac12+\varepsilon$ fraction of inputs, matching
      bounds of $\varepsilon^2n^{3/2}$ gates up to polylogarithmic factors
      (for some $\varepsilon$, see below):
      an explicit circuit with $O(\varepsilon^2n^{3/2}\sqrt{\log(1/\varepsilon)})$
      gates, and a lower bound of $\Omega(\varepsilon^2n^{3/2}/\log^2n)$
      (previously $\varepsilon^3$ and $\log^3n$);
\item worst-case lower bounds of $\Omega(n^{3/2}/\log^2 n)$ gates and
      $\Omega(n^{5/2}/\log^{5/2}n)$ wires (previously $\log^3 n$ and
      $\log^{7/2}n$ in the denominators).
\end{itemize}
In our circuits, and in that of Kane and Williams, each bottom gate's weights
have a single sign within every address block. We show that such circuits
for Andreev's function need $\Omega(n^2/\log^3n)$ gates. Within this class our
exact circuit is therefore optimal up to polylogarithmic factors, and beating
$n^2$ by more than that requires mixed signs.

All numbered results are formalized in Lean; the asymptotic consequences
are derived from explicit bounds in the parameters.
\end{abstract}

%===================================================================
\section{Setting and results}

Lower bounds for depth-two threshold circuits are a classical frontier of
circuit complexity: for explicit functions in $\mathsf P$, nothing beyond
roughly $n^{3/2}$ gates is known. The first superlinear gate lower bound is
due to Kane and Williams~\cite{KW}, whose hard function is Andreev's function: look up a bit
in a table, at an address given by parities. They proved
$\Omega(n^{3/2}/\log^3n)$ gates, gave an $O(n^3/\log n)$ circuit, and asked for
asymptotically tight bounds. This article narrows the gap from both sides,
settles the approximate version up to polylogarithmic factors, and locates
what remains of the exact gap: if every bottom gate has weights of a single
sign inside each address block, then $\Omega(n^2/\log^3n)$ gates are
necessary, so a substantially subquadratic circuit must have bottom gates
with weights of both signs inside some address block.

\paragraph*{Andreev's function.}
Fix integers $k,m\ge1$. The input consists of
\begin{itemize}
\item a \emph{table} $x$ of $2^k$ bits $x_c$, one for each address
      $c\in\bits^k$, and
\item $k$ \emph{address blocks} of $m$ bits each; block $i$ consists of the bits
      $a_{i,1},\dots,a_{i,m}$.
\end{itemize}
The function is
\[
A(x,a)=x_b,\qquad\text{where }b_i=\left(\sum_{j=1}^m a_{i,j}\right)\bmod 2\quad\text{for }i=1,\dots,k.
\]
Write $N=2^k$ for the length of the table, so the input length is $n=N+km$.
In the \emph{balanced case} the table and the address part have equal
length, $km=N$, so $n=2N$ and $k=\log_2(n/2)$. Exact balance requires
$k$ to be a power of two. Asymptotic statements refer to this balanced
sequence; statements for general $k,m$ do not require this assumption.
We use $\log$ for the base-two logarithm and $\ln$ for the natural logarithm.

\paragraph*{Threshold circuits.}
For a statement $P$, the bracket $\ind{P}$ is $1$ if $P$ holds and $0$
otherwise. A \emph{threshold gate} with inputs $z_j\in\bits$, real weights
$\omega_j$ and threshold $\tau$ computes $\ind{\sum_j\omega_jz_j\ge\tau}$. In a \emph{depth-two threshold
circuit} ($\THR\circ\THR$) some \emph{bottom} gates read input bits, and a
single \emph{top} gate reads the bottom gates and input bits. Its
\emph{size} is its total number of gates, top gate included; its
\emph{wire count} is the number of connections with nonzero weight.
For $0<\varepsilon\le1/2$, a circuit \emph{$\varepsilon$-agrees} with a
Boolean function if it outputs the function's value on at least a
$\tfrac12+\varepsilon$ fraction of all inputs. Agreement fractions are
measured under the uniform distribution on inputs.

\paragraph*{Results and the state of the art.}
Stronger lower bounds against $\THR\circ\THR$ are known for functions in the
much larger class $\mathsf E^{\mathsf{NP}}$; they are proved by the
algorithmic method~\cite{Tam,ACW,CTW}. For functions in $\mathsf P$, which
includes Andreev's function, the best known gate lower bounds are of
order $n^{3/2}/\mathrm{polylog}\,n$~\cite{KW,KKL}, as also noted by Chen,
Tal and Wang~\cite{CTW}. Our bounds improve the explicit polylogarithmic factors and
the dependence on $\varepsilon$.
\begin{center}
\small
\setlength{\tabcolsep}{5pt}
\begin{tabular}{@{}lllll@{}}
\toprule
Function & Bound & Best previously known & This article & \\
\midrule
\multicolumn{5}{@{}l}{\emph{Functions in $\mathsf P$}}\\
Andreev, exact & gates, upper & $O(n^3/\log n)$ \cite{KW} & $(5/16+o(1))n^2$ & Thm.~\ref{thm:exact}\\
Andreev, exact & gates, lower & \begin{tabular}[t]{@{}l@{}}$\Omega(n^{3/2}/\log^3n)$ \cite{KW}\\ $n^{3/2}/(\log n)^{O(1)}$ \cite{KKL}$^\dagger$\end{tabular} & $\Omega(n^{3/2}/\log^2n)$ & Thm.~\ref{thm:lower}\\
Andreev, exact & wires, lower & $n^{5/2}/(\log n)^{O(1)}$ \cite{KKL}$^\dagger$ & $\Omega(n^{5/2}/\log^{5/2}n)$ & Thm.~\ref{thm:lower}\\
Andreev, exact$^\ddagger$ & gates, lower & none stated & $\Omega(n^2/\log^3n)$ & Thm.~\ref{thm:coherent}\\
Andreev, $\varepsilon$-agree & gates, upper & none stated & $O(\varepsilon^2n^{3/2}\sqrt{\log(1/\varepsilon)})$ & Thm.~\ref{thm:avgupper}\\
Andreev, $\varepsilon$-agree & gates, lower & $\Omega(\varepsilon^3n^{3/2}/\log^3n)$ \cite{KW} & $\Omega(\varepsilon^2n^{3/2}/\log^2n)$ & Thm.~\ref{thm:lower}\\
Andreev, $\varepsilon$-agree & wires, lower & $\Omega(\varepsilon^3n^{5/2}/\log^{7/2}n)$ \cite{KW} & $\Omega(\varepsilon^2n^{5/2}/\log^{5/2}n)$ & Thm.~\ref{thm:lower}\\
\midrule
\multicolumn{5}{@{}l}{\emph{Functions in $\mathsf{E}^{\mathsf{NP}}$}}\\
some $F$, exact & gates, lower & \begin{tabular}[t]{@{}l@{}}$n^{2-\delta}$ \cite{Tam,ACW}\\ $n^{5/2-\delta}$ \cite{CTW}\end{tabular} & -- & \\
some $F$, $\varepsilon$-agree & wires, lower & $n^3/\log^{12}n$ \cite{Nag} (unreviewed) & -- & \\
\bottomrule
\end{tabular}
\end{center}
Notes on the table:
\begin{itemize}
\item All rows concern $\THR\circ\THR$ circuits.
      The Andreev rows are asymptotic in the balanced case; its exact upper
      bound has the explicit finite version \eqref{eq:bestcount}.
\item The $\varepsilon$-agreement rows for Andreev's function hold for
      $\varepsilon\ge3\sqrt2(\log_2n+1)/\sqrt n$. Below this, a single
      majority gate over the table already agrees on a
      $\tfrac12+\Omega(1/\sqrt n)$ fraction. The upper bound is stated for
      $\varepsilon\le\tfrac18$. For $\tfrac18<\varepsilon<\tfrac12$, use
      Part~1 of Theorem~\ref{thm:avgupper} with $\delta=\tfrac12-\varepsilon$;
      at $\varepsilon=\tfrac12$, use the exact construction.
\item $^\dagger$For a variant of Andreev's function whose table is four times
      as long as the address part.
\item $^\ddagger$For circuits whose bottom gates are \emph{sign-coherent}:
      inside every address block, the weights of a gate are all $\ge0$ or
      all $\le0$ (Section~\ref{sec:coherent}). The circuit of
      Theorem~\ref{thm:exact} is of this kind.
\end{itemize}

The remaining gap for unrestricted Andreev circuits concerns exact computation: between
$n^{3/2}/\log^2n$ and $(5/16+o(1))n^2$ gates. For circuits with sign-coherent
bottom gates the gap is only polylogarithmic. For any fixed $\delta>0$, a circuit with $n^{2-\delta}$
gates therefore needs, for sufficiently large $n$, bottom gates whose
weights have mixed signs inside a block. The two
polynomial-size constructions we know of avoid such gates: besides ours, the
circuit of Kane and Williams~\cite[Thm.~1.2]{KW} builds each of its parities
from gates that threshold the plain sum of the parity's
variables~\cite[Prop.~2.1]{HMPST}, so it is sign-coherent as well.

\paragraph*{Related work.}
Andreev's function comes from formula complexity: for De Morgan formulas,
shrinkage gives it a lower bound of $n^{3-o(1)}$, the best known for an
explicit function~\cite{And87,Has98,Tal14}. For depth-two threshold circuits,
exponential lower bounds for inner product hold when the top gate has small
weights~\cite{HMPST,Nis93}. With arbitrary weights at both
levels, the bounds of Kane and Williams~\cite{KW} were the first superlinear
gate and superquadratic wire bounds; Kabanets, Kane and Lu~\cite{KKL} and Chen,
Santhanam and Srinivasan~\cite{CSS} refine their restriction method and use it
for satisfiability algorithms, and the algorithmic method gives stronger
bounds for functions in $\mathsf E^{\mathsf{NP}}$~\cite{Tam,ACW,CTW}.
Sign-coherent gates are unate in each address block; unateness also drives
the parity lower bounds of Impagliazzo, Paturi and Saks~\cite{IPS}.

\paragraph*{How to read this article.}
Section~\ref{sec:tech} lists the new techniques and where they are
used. Sections~\ref{sec:upper}--\ref{sec:avgupper} (upper bounds),
Section~\ref{sec:lower} (lower bounds) and Section~\ref{sec:coherent}
(sign-coherent circuits) can be read independently. Section~\ref{sec:open} lists open questions.

%===================================================================
\section{New techniques}\label{sec:tech}

The proofs combine standard tools, namely Fourier expansion of the lookup,
threshold staircases for functions of a linear form (as for symmetric
functions~\cite[Prop.~2.1]{HMPST}),
and the random restrictions of Kane and Williams~\cite{KW}, with the following
new ingredients.

\begin{itemize}
\item \textbf{Removing the parity factor.} Used for the exact upper
      bound (Lemma~\ref{lem:parity}, Proposition~\ref{prop:basic},
      Theorem~\ref{thm:exact}). The Fourier
      expansion writes the lookup as a sum of products ``table coefficient
      $\times$ address parity''. For $\varphi(u)\coloneqq(-1)^uu$ and even $X$,
      $X(-1)^t=\varphi(X+t)-\varphi(t)$, so every
      product becomes the difference of two functions of affine forms,
      each implemented by a threshold staircase. Kane and Williams build each
      of about $n^2/\log n$ parities
      separately. Collecting the terms that share the same address parity
      $(-1)^{S\cdot b}$ leaves only $N$ terms, each with an integer
      coefficient $X_S$, and saves a factor $n/\log n$. Sharing quadratic corrections across characters and
      absorbing errors at a few special addresses gives the constant $5/16$.
\item \textbf{Odd error functions are orthogonal.} Used for the upper bounds for
      approximate computation (Lemma~\ref{lem:odd},
      Theorem~\ref{thm:avgupper}). For each Fourier coefficient, we pair two
      reflected forms and truncate them to the same symmetric window.
      Their combined error is an odd function of that coefficient. A symmetry of the uniform table flips one coefficient and
      fixes another, so the errors of different characters are orthogonal
      although the coefficients are not independent. Their squares add, and
      no union bound over characters is needed. Keeping only the characters of
      the first $r$ address bits gives a majority vote over a subcube of the
      table and the $\varepsilon^2$ tradeoff.
\item \textbf{Fixed-weight counting.} Used for both lower bounds, the
      general one and the one for sign-coherent circuits
      (Lemma~\ref{lem:count}, Theorems~\ref{thm:lower} and~\ref{thm:coherent}). After a restriction, a
      surviving gate keeps its original weights on the live bits; only its
      threshold moves. This fact is standard~\cite{KW,CSS,CTW}, but we use it
      for counting: a survivor costs $O(k+\log s)$ bits instead of the
      $\Theta(k^2)$ bits of an arbitrary threshold function~\cite{KW,KKL,CSS}.
      This saves one logarithm. It works because the circuit is fixed before
      the random table is drawn.
\item \textbf{A union bound uniform in the number of survivors.} Used
      for the lower bound for approximate computation (Lemma~\ref{lem:avg},
      Theorem~\ref{thm:lower}). Instead of bounding the
      survivors by Markov's inequality, we bound the correlation for all
      survivor counts at once and average with Cauchy--Schwarz. This improves
      $\varepsilon^3$ to $\varepsilon^2$.
\item \textbf{The cube restriction.} Used for the lower bound for
      sign-coherent circuits (Lemmas~\ref{lem:entropy} and~\ref{lem:diag},
      Theorem~\ref{thm:coherent}). Instead of fixing the other address bits at
      random, fix a prefix of every block and let its length run through all
      values. On every resulting unit cube of block counts the circuit still
      reveals the whole table, while a sign-coherent gate moves monotonically
      through the grid and changes value at most once along each diagonal
      line, a classical fact for unate functions~\cite[Lemma~2.1]{IPS}. So it
      is non-constant on at most a $k/m$ fraction of the cubes, compared
      with the $O(\sqrt n/m)$ survival bound for random restrictions.
\end{itemize}

%===================================================================
\section{Exact quadratic circuits}\label{sec:upper}

\paragraph*{Notation for Sections~\ref{sec:upper} and~\ref{sec:avgupper}.}
Let $|a_i|$ be the number of ones in block $i$, so the address is $b$
with $b_i=|a_i|\bmod2$.
For $S\subseteq\{1,\dots,k\}$ and $c\in\bits^k$ write $S\cdot c=\sum_{i\in S}c_i$, and
\[
X_S(x)\coloneqq\sum_{c\in\bits^k}(-1)^{S\cdot c+x_c},\qquad
t_S(a)\coloneqq\sum_{i\in S}|a_i| .
\]
$X_S(x)$ is $N$ times the Fourier coefficient at $S$ of $c\mapsto(-1)^{x_c}$;
it is even and lies in $[-N,N]$. $t_S(a)$ is the number of ones in the blocks in
$S$, and lies in $[0,|S|m]$. We drop the arguments when they are clear.
We write $[r]=\{1,\dots,r\}$, and $c_{[r]}$ for the first $r$ coordinates of $c$.

\paragraph*{Plan.}
Testing the sign of the score $N(-1)^{x_b}$ computes Andreev's function.
Lemma~\ref{lem:lookup} writes it as a sum of products,
Lemma~\ref{lem:parity} removes the products, and Lemma~\ref{lem:ladder}
turns the result into threshold gates. We then refine this to a smaller score
with the same sign, proving the $5/16$ bound.

\begin{lemma}[Lookup as a sum of products]\label{lem:lookup}
For every table $x$ and address input $a$,
\[
\sum_{S\subseteq\{1,\dots,k\}}X_S(x)\,(-1)^{t_S(a)}=N\,(-1)^{x_b}.
\]
\end{lemma}
\begin{proof}
Since $|a_i|\equiv b_i\pmod2$, we have $(-1)^{t_S}=(-1)^{S\cdot b}$.
Exchanging the two sums, the left side is
\[
\sum_c(-1)^{x_c}\sum_{S\subseteq\{1,\dots,k\}}(-1)^{S\cdot c+S\cdot b}
=\sum_c(-1)^{x_c}\sum_{S\subseteq\{1,\dots,k\}}\ \prod_{i\in S}(-1)^{c_i+b_i}.
\]
Choosing $S$ means choosing, for each $i$, whether $i\in S$, so the inner sum
factors as $\prod_{i=1}^k\bigl(1+(-1)^{c_i+b_i}\bigr)$. This product is $N$ for
$c=b$ and $0$ otherwise.
\end{proof}

\begin{lemma}[Removing the parity factor]\label{lem:parity}
Define $\varphi$ on the integers by
\[
\varphi(\theta)\coloneqq(-1)^\theta\,\theta .
\]
If $X$ is an even integer and $t$ is an integer, then
$X(-1)^t=\varphi(X+t)-\varphi(t)$.
\end{lemma}
\begin{proof}
Because $X$ is even, $X+t$ has the parity of $t$, so $\varphi(X+t)=(-1)^t(X+t)$.
\end{proof}

\begin{lemma}[Threshold staircase]\label{lem:ladder}
Let $L$ be an integer-valued quantity with $\ell\le L\le h$, where
$\ell,h$ are integers, and let $g:\mathbb Z\to\mathbb Z$. Then
\[
g(L)=g(\ell)+\sum_{\theta=\ell+1}^{h}\bigl(g(\theta)-g(\theta-1)\bigr)\ind{L\ge\theta}.
\]
\end{lemma}
\begin{proof}
The brackets equal $1$ exactly for $\theta\le L$, and the sum telescopes.
\end{proof}

When $L$ is an affine form in the input bits, each $\ind{L\ge\theta}$ is one
threshold gate, so $g(L)$ costs one gate per value of $\theta$, i.e.\ $h-\ell$
gates. Both $t_S$ and $X_S+t_S$ are such forms, since
$(-1)^{x_c}=1-2x_c$:
\begin{equation}\label{eq:form}
X_S+t_S=\sum_c(-1)^{S\cdot c}-2\sum_c(-1)^{S\cdot c}x_c+\sum_{i\in S}\sum_{j=1}^m a_{i,j}.
\end{equation}
We now combine Lemmas~\ref{lem:lookup}--\ref{lem:ladder} into a quadratic
circuit for Andreev's function. Theorem~\ref{thm:exact} below refines it to
$(5/16+o(1))n^2$ gates, with the proof in Appendix~\ref{app:exact}.

\begin{proposition}[A first quadratic construction]\label{prop:basic}
For all $k,m\ge1$, Andreev's function has an exact depth-two threshold
circuit with $2N^2+kmN$ bottom gates and one top gate. In the balanced case
this is $\tfrac34n^2+1$ gates. Bottom weights lie in $\{0,\pm2\}$ on table
bits and in $\{0,1\}$ on address bits, and
top weights are integers of absolute value at most $2n-1$.
\end{proposition}
\begin{proof}
By Lemmas~\ref{lem:lookup} and~\ref{lem:parity},
$N(-1)^{x_b}=\sum_S\left(\varphi(X_S+t_S)-\varphi(t_S)\right)$.
Apply Lemma~\ref{lem:ladder} to every term, with $X_S+t_S$ ranging over
$[-N,\,N+|S|m]$ and $t_S$ over $[0,|S|m]$. Put
$d_\theta\coloneqq\varphi(\theta)-\varphi(\theta-1)=(-1)^\theta(2\theta-1)$.
Since $N$ is even, $\varphi(-N)=-N$, and $\varphi(0)=0$, so
\[
N(-1)^{x_b}=-N^2+\sum_S\Biggl(\sum_{\theta=1-N}^{N+|S|m}d_\theta\,\ind{X_S+t_S\ge\theta}
-\sum_{\theta=1}^{|S|m}d_\theta\,\ind{t_S\ge\theta}\Biggr).
\]
The bottom gates are the indicators $\ind{X_S+t_S\ge\theta}$ and
$\ind{t_S\ge\theta}$, threshold gates by~\eqref{eq:form}. The top gate
outputs $1$ exactly when the double sum is less than $N^2$, that is, when
$x_b=1$. The number of bottom gates is $\sum_S(2N+2|S|m)=
2N^2+2m\sum_S|S|=2N^2+kmN$, where $\sum_S|S|=kN/2$ because
each block $i\in\{1,\dots,k\}$ lies in half of the $N$ sets $S$.
In the balanced case $km=N$ and $n=2N$, so this is $3N^2=\tfrac34n^2$.
The top weights are $\pm d_\theta$. Since $-N<\theta\le N+km=n$, their
absolute value is at most $2n-1$.
\end{proof}
Since the bottom weights are small integers, the circuit is even a threshold
of majority gates ($\THR\circ\mathsf{MAJ}$, by repeating inputs). For
comparison, Kane and Williams also give a $\mathsf{MOD}_3\circ\mathsf{MOD}_2$
circuit of size $O(n^2)$~\cite[Thm.~1.2]{KW}, in a different basis.

% Copyright (c) 2026 Christian Reitwiessner.
% Authors: Christian Reitwiessner
\subsection*{A sharper score: the \texorpdfstring{$5/16$}{5/16} bound}

\begin{theorem}\label{thm:exact}
Put $\alpha=\lfloor(m-1)/2\rfloor$. There is an exact depth-two threshold
circuit for $A$ with at most
\begin{equation}\label{eq:bestcount}
\frac{N(N-1)}2+\frac{kN(m+\alpha)}2
 +2(N-1)+2\alpha k+3km+\left\lceil\frac{km}{2}\right\rceil
\end{equation}
bottom gates and one top gate. Its weights are integers of magnitude
polynomial in $n$. In the balanced case its total size is
\[
\frac5{16}n^2-\frac n4\log_2(n/2)+O(n).
\]
\end{theorem}

\paragraph*{Idea of the proof.}
Since the top gate only tests the sign of $N(-1)^{x_b}$, any score that is
at least $N$ when $x_b=0$ and at most $-N$ when $x_b=1$ serves. With
$Y_S=\sum_{c:\,S\cdot c\text{ odd}}x_c$ we have $X_S=4Y_S-2\sum_cx_c$ for
$S\ne\varnothing$, so replacing $X_S$ by $4Y_S$ in Lemma~\ref{lem:lookup}
changes the sum only by an affine term and a multiple of $\ind{b=0}$, and
$u_S=2Y_S+t_S$ reduces the width of the table-dependent part of the range
from $2N$ to $N$. Further, $4Y_S(-1)^{t_S}$
equals $-8\lceil u_S/2\rceil^2$ up to a polynomial in $Y_S$, $t_S$ and
$\lceil t_S/2\rceil$, and the square changes only at odd $u_S$, so its
staircase needs only odd thresholds. Summed over $S$, the polynomials collapse
into $O(n)$ gates, up to errors at $b=0$ and at the $k$ addresses $e_i$ of
weight one. A term $-Cx_0\ind{b=0}$ with $C$ large, built from the forms
$t_S$ with $x_0$ attached by a larger weight, repairs $b=0$ at about $kmN/4$
gates. At $b=e_i$ the error is $-N|a_i|$
with $|a_i|\le2\alpha+1$; a shift of the staircase coefficients handles
$x_{e_i}=0$, and weight $1+\alpha$ on $x_{e_i}$ inside $Y_S$ handles
$x_{e_i}=1$. The full proof is in Appendix~\ref{app:exact}.

The sharper construction remains sign-coherent: every bottom gate sees a block
through its sum with a nonnegative coefficient. Mixed-sign weights
are therefore not needed for this bound. Section~\ref{sec:coherent} shows
that improving the quadratic exponent requires gates that are not
sign-coherent.


%===================================================================
% Copyright (c) 2026 Christian Reitwiessner.
% Authors: Christian Reitwiessner
\section{Agreeing on most inputs with fewer gates}\label{sec:avgupper}

Under uniform inputs every form $X_S\pm t_S$ concentrates in a window of
width $O(\sqrt n)$, so each staircase can be cut to its window. The question
is how the $N$ cut-off errors add up. They are not independent, but with a
symmetric truncation they are orthogonal.

\begin{lemma}[Odd error functions are orthogonal]\label{lem:odd}
Let the table be uniformly random. For distinct sets $S,T\subseteq\{1,\dots,k\}$ and functions
$g,h:\mathbb Z\to\mathbb R$ with $g$ an odd function, that is, $g(-y)=-g(y)$ for all $y$, we have
$\mathbb E\bigl[g(X_S)\,h(X_T)\bigr]=0$.
\end{lemma}
\begin{proof}
Since $S\ne T$, pick $u\in\bits^k$ with $(S\mathbin\triangle T)\cdot u$ odd,
and map the table $x$ to $x'$ with $(-1)^{x'_c}=(-1)^{T\cdot u}(-1)^{x_{c\oplus u}}$.
This is a bijection of tables, so $x'$ is uniform too, and
\begin{align*}
X_S(x')
 &=(-1)^{T\cdot u}\sum_{c}(-1)^{S\cdot(c\oplus u)+x_c}
  =(-1)^{S\cdot u+T\cdot u}\sum_{c}(-1)^{S\cdot c+x_c}\\
 &=(-1)^{(S\mathbin\triangle T)\cdot u}X_S(x)=-X_S(x),
\end{align*}
and in the same way $X_T(x')=X_T(x)$. Since $x'$ is uniform,
\[
\mathbb E_x[g(X_S(x))h(X_T(x))]=\mathbb E_{x'}[g(X_S(x'))h(X_T(x'))]
=\mathbb E_x[g(-X_S(x))h(X_T(x))].
\]
As $g$ is an odd function, the right side is the negative of the left side, so both are zero.
\end{proof}
$X_S$ and $X_T$ are not independent; the lemma only uses this
pairwise sign symmetry.

\begin{theorem}\label{thm:avgupper}
There are explicit depth-two threshold circuits for Andreev's function $A$
with $m=N/k$ (the balanced case) as follows:
\begin{enumerate}
\item For $0<\delta\le1/2$, a circuit with $O\bigl(n^{3/2}\sqrt{\log(1/\delta)}\bigr)$
      gates that agrees with $A$ on at least a $1-\delta$ fraction of all inputs.
\item For $N^{-1/2}\le\varepsilon\le1/8$, a circuit with
      $O\bigl(\varepsilon^2n^{3/2}\sqrt{\log(1/\varepsilon)}\bigr)$ gates and
      $O\bigl(\varepsilon^2n^{5/2}\sqrt{\log(1/\varepsilon)}\bigr)$ wires that
      $\varepsilon$-agrees with $A$, i.e.\ agrees on at least a $1/2+\varepsilon$ fraction.
\end{enumerate}
All these circuits have sign-coherent bottom gates, i.e.\
inside every address block the weights of a gate are all $\ge0$ or all
$\le0$.
\end{theorem}

\begin{proof}
\emph{A majority vote with fewer characters.}
Fix $0\le r\le k$ and put $q=2^r$. The proof of Lemma~\ref{lem:lookup} with
the sum restricted to $S\subseteq[r]$ gives
\[
P_*\coloneqq2\sum_{S\subseteq[r]}X_S(-1)^{t_S}
 =2q\sum_{c:\,c_{[r]}=b_{[r]}}(-1)^{x_c}.
\]
So $P_*$ is a vote among $L=N/q$ table bits, one of them $x_b$. For $q=N$
we have $P_*=\pm2N$ and output $\ind{-P_*\ge0}$; for $q<N$ we output the
majority $\ind{-P_*\ge2q}$, with ties going to $0$.

For $q<N$, fix the address input; the $L$ voters are then independent
uniform bits, and $L$ is even. In the case where exactly $L/2$ of the other
$L-1$ voters are $1$, the vote follows $x_b$ and is correct. Otherwise it does
not depend on $x_b$, so it is correct with probability $1/2$. Hence its agreement is
\[
\frac12+\frac12\cdot2^{-(L-1)}\binom{L-1}{L/2}=\frac12+2^{-L-1}\binom{L}{L/2}
\ge\frac12+\frac1{2\sqrt{2L}}.
\]

\emph{Symmetric truncation.}
For nonempty $S\subseteq[r]$ center the address count,
$D_S=t_S-\lfloor m|S|/2\rfloor$, and put
$\eta_S=(-1)^{\lfloor m|S|/2\rfloor}$. Lemma~\ref{lem:parity}, applied with
$t=D_S$ and $t=-D_S$, gives
\[
\eta_S\bigl(\varphi(X_S+D_S)+\varphi(X_S-D_S)\bigr)=2X_S(-1)^{t_S}.
\]
For an integer radius $R\ge1$, let $P_R$ be obtained from
$P_*=2X_\varnothing+\sum_{S\ne\varnothing}\eta_S\bigl(\varphi(X_S+D_S)+\varphi(X_S-D_S)\bigr)$
by replacing every $\varphi(u)$ by its truncation $\varphi(u)\ind{|u|\le R}$
(the affine term $2X_\varnothing$ stays). Each truncated staircase has at most
$2R+2$ steps, so $P_R$ uses at most $4(q-1)(R+1)$ bottom gates. Their
address weights are $+1$ on the blocks of $S$ in one form and $-1$ in the
other, so they are sign-coherent.

\emph{Cancellation of cross terms.}
Let $e_S$ be the term of $S$ in $P_R$ minus its term in $P_*$; the affine term
cancels, so $P_R-P_*=\sum_{S\ne\varnothing}e_S$. Fix the address input. Then
$D_S$ and $\eta_S$ are constants and
\[
e_S=-\eta_S\bigl(\varphi(X_S+D_S)\ind{|X_S+D_S|>R}+\varphi(X_S-D_S)\ind{|X_S-D_S|>R}\bigr)
\eqqcolon g_S(X_S)
\]
is a function of $X_S$ alone. The function $g_S$ is odd because
$u\mapsto\varphi(u)\ind{|u|>R}$ is an odd function: $\varphi(-u)=-\varphi(u)$
and the cutoff depends only on $|u|$. Replacing $y$ by $-y$ sends $y+D_S$
to $-(y-D_S)$ and $y-D_S$ to $-(y+D_S)$, so the two summands are exchanged
and negated, giving $g_S(-y)=-g_S(y)$. The table stays uniform when the address input is fixed,
so by Lemma~\ref{lem:odd} with $g=g_S$ and $h=g_T$ we get
$\mathbb E[e_Se_T]=0$ for all $S\ne T$, first for each fixed address and then
on average. Therefore
\[
\mathbb E(P_R-P_*)^2=\sum_{S\ne\varnothing}\mathbb E\,e_S^2
+\sum_{S\ne T}\mathbb E[e_Se_T]=\sum_{S\ne\varnothing}\mathbb E\,e_S^2 .
\]

\emph{Diagonal terms.}
Here the expectations are over both the table and the address bits.
For $U=X_S\pm D_S$ we have $|\mathbb EU|\le1/2$, and $U-\mathbb EU$ is a sum
of $N$ independent uniform signs and $m|S|$ independent uniform
variables in $\{-1/2,1/2\}$, so Hoeffding's inequality
gives $\Pr\bigl[|U-\mathbb EU|\ge\lambda\bigr]\le2\exp(-\lambda^2/(2v))$ with
$v=N+km/4$. As $|U|>R$ implies $|U-\mathbb EU|>R/2$ and
$U^2\le4(U-\mathbb EU)^2$, integrating this tail from $R/2$ gives
$\mathbb E\bigl[U^2\ind{|U|>R}\bigr]\le(2R^2+16v)e^{-R^2/(8v)}
\le32v\exp\bigl(-R^2/(16v)\bigr)$, and $e_S^2\le2\sum_\pm U^2\ind{|U|>R}$
gives
\[
\mathbb E(P_R-P_*)^2=\sum_{S\ne\varnothing}\mathbb E\,e_S^2
\le128(q-1)v\exp\bigl(-R^2/(16v)\bigr).
\]
For $q<N$, the sum $P_*/(2q)$ contains an even number $L$ of signs, so
$P_*$ is a multiple of $4q$ and lies at distance at least $2q$ from the
threshold $-2q$. For $q=N$, the two scores $\pm2N$ have the same margin
from the threshold $0$. Thus the truncated circuit
differs from the vote only if $|P_R-P_*|\ge2q$. By Markov's inequality this
has probability at most $(32v/q)\exp(-R^2/(16v))$, which is at most $\delta$
for
\[
R\coloneqq\left\lceil4\sqrt{v\ln\frac{32v}{q\delta}}\right\rceil .
\]

\emph{The two parameter choices.}
For Part~1 take $q=N$, so the vote is $A$ itself; at balance $v=O(N)$ and
$R=O(\sqrt{N\log(1/\delta)})$. For Part~2 take the least power of two
$q\ge32\varepsilon^2N$ and $\delta=\varepsilon$. Then $32\le q\le N/2$
because $N^{-1/2}\le\varepsilon\le1/8$, and $q<64\varepsilon^2N$; the vote
agrees on at least a $1/2+\sqrt{q/(8N)}\ge1/2+2\varepsilon$ fraction, and
truncation loses at most $\varepsilon$. Every bottom gate has at most $n$
input wires and one wire to the top
gate, which also has at most $n$ direct input wires. This gives the stated
wire bound.
\end{proof}


%===================================================================
\section{Lower bounds}\label{sec:lower}

This section proves lower bounds on the gates and wires of an arbitrary
depth-two threshold circuit that agrees with Andreev's function on at least a
$1/2+\varepsilon$ fraction of inputs; $\varepsilon=1/2$ means exact computation.
For gates, the bound matches Theorem~\ref{thm:avgupper} up to
polylogarithmic factors. The argument is the random restriction method of
Kane and Williams~\cite{KW} with two changes taken from
Section~\ref{sec:tech}: a sharper count of restricted circuits saves a
logarithm, and a union bound over all survivor counts improves
$\varepsilon^3$ to $\varepsilon^2$.

\begin{theorem}\label{thm:lower}
Let $C$ be a depth-two threshold circuit with $s$
gates and $w$ wires that $\varepsilon$-agrees with Andreev's function, where
$\varepsilon\ge3(k+2)/\sqrt N$. Then, in the balanced case ($km=N$),
\[
s=\Omega\bigl(\varepsilon^2n^{3/2}/\log^2n\bigr)
\quad\text{and}\quad
w=\Omega\bigl(\varepsilon^2n^{5/2}/\log^{5/2}n\bigr).
\]
Taking $\varepsilon=\tfrac12$,
allowed when $N\ge36(k+2)^2$, gives the worst-case bounds.
\end{theorem}

\paragraph*{Plan.}
We follow Kane and Williams~\cite{KW}. Fix all but one bit per block at
random. The restricted target function is a lookup in a \emph{random} table
indexed by the $k$ remaining bits, and only a few bottom gates still depend
on those bits (Lemma~\ref{lem:surv}). There are few such restricted circuits
(Lemma~\ref{lem:count}), and with high probability none correlates well with a random
table (Lemma~\ref{lem:avg}). Averaging over all restrictions proves the
theorem.

\paragraph*{Restrictions.}
The lemmas below hold for all $k,m\ge1$, not only in the balanced case.
Throughout, $C$ is a depth-two threshold circuit of size $s$ with $w$ wires.
A \emph{restriction} consists of a choice $\sigma$ of one \emph{live}
position in every block, and a \emph{fixing}: values for all other input
bits, table included. For a uniformly random restriction, the live position
in each block is
chosen uniformly and independently, and all fixed bits are independent
uniform bits, independent of these choices. Write $u\in\bits^k$ for the
live bits and
$o_i$ for the parity of the $m-1$ fixed bits of block $i$ (the
\emph{offset}). Then the address is $b=u\oplus o$, and the restricted
function is $u\mapsto x_{u\oplus o}$. A bottom gate \emph{survives} if it
still depends on $u$. Let $U$ be the number of surviving gates.

\begin{lemma}[Few survivors]\label{lem:surv}
For a uniformly random restriction, $\mathbb E[U]\le4s\sqrt n/m$. The
expected number of bottom gates that depend on at least two live bits is at
most $4w/m^{3/2}$. Both bounds also hold for every fixed table, averaging
over the live
positions and the fixed address bits.
\end{lemma}
These are the random restriction lemmas of Kane and
Williams~\cite[Lemmas~1.1 and~3.1]{KW}, with explicit constants.
Appendix~\ref{app:surv} proves both bounds; the proof includes small repairs
to their argument (ties among weights, zero weights, and a strict
versus non-strict inequality). Chen, Tal and Wang give an alternative proof
of the first lemma, up to constants~\cite[App.~A.3]{CTW}; Kabanets, Kane and
Lu~\cite[Lemma~6.1]{KKL} and Chen, Santhanam and
Srinivasan~\cite[Lemma~9.1]{CSS} prove analogues with polylogarithmic and
$n^{o(1)}$ losses.

\begin{lemma}[Roychowdhury, Siu, Orlitsky~\cite{RSO}; see {\cite[Thm.~2.3]{KW}}]\label{fact:rso}
For any fixed Boolean functions $f_1,\dots,f_q:\Omega\to\bits$ on a finite
set $\Omega$, at most
$(|\Omega|+1)^{q+1}$ functions on $\Omega$ have the form
$\ind{\sum_j\omega_jf_j\ge\tau}$.
\end{lemma}

\begin{lemma}[Counting restricted circuits]\label{lem:count}
Fix $C$ and $\sigma$. Every fixing turns $C$ into a function of the live
bits $u\in\bits^k$. For an integer $q\ge0$, let $\cH_q$ be the set of
functions that arise in this way from fixings with at most $q$ survivors. Then
\[
|\cH_q|\le\bigl((2^k+1)(s+1)\bigr)^q\,(2^k+1)^{q+k+1}.
\]
\end{lemma}
\begin{proof}
Non-surviving bottom gates are constants, and are absorbed into the top
threshold. A surviving gate is $\ind{\omega\cdot u\ge\tau}$, where $\omega$ is
its original weight vector restricted to the live positions; only $\tau$
depends on the fixing. For a fixed $\omega$, these functions are nested as
$\tau$ grows, so there are at most $2^k+1$ of them. Pad the list of survivors
with constant functions to length exactly $q$. Then each entry of the list is
one of at most $(2^k+1)(s+1)$ functions: choose one of the at most $s$ gates
or ``constant'', then the threshold. The top gate is a threshold function of
the $q$ entries and the $k$ live bits. By Lemma~\ref{fact:rso} with
$|\Omega|=2^k$, there are at most $(2^k+1)^{q+k+1}$ such functions for each
list.
\end{proof}
Kane and Williams~\cite{KW}, and similarly Kabanets, Kane and Lu~\cite{KKL}
and Chen, Santhanam and Srinivasan~\cite{CSS}, count all threshold
functions of the live bits at this step, which costs about $k^2$ bits per
survivor instead of $O(k+\log s)$.

\begin{lemma}[Agreement bound for all survivor counts at once]\label{lem:avg}
Let $\alpha=\sqrt{N(k+2+\ln(s+1))}$ and $\beta=(k+2)\sqrt N$. Then $C$
agrees with Andreev's function on at most a
\[
\tfrac12+\frac{\alpha\sqrt{\mathbb E[U]}+\beta}{N}+2^{-(k+6)}
\]
fraction of inputs.
\end{lemma}
\begin{proof}
Fix $\sigma$ and the fixed address bits; the table remains uniformly random.
For one function $f$ of $u$ and $\Delta\ge0$, Hoeffding's inequality gives
\[
\Pr\bigl[f\text{ agrees with }u\mapsto x_{u\oplus o}\text{ on}\ge N/2+\Delta\text{ points}\bigr]
\le e^{-2\Delta^2/N}.
\]
Put $\Delta_q=\alpha\sqrt q+\beta$. Call the table \emph{bad} if, for some
integer $0\le q<N$, a function $f\in\cH_q$ agrees with
$u\mapsto x_{u\oplus o}$ on more than $N/2+\Delta_q$ points. Using $\Delta_q^2\ge\alpha^2q+\beta^2$,
$\ln(2^k+1)\le k+1$ and Lemma~\ref{lem:count},
\begin{align*}
\Pr[\text{bad}]&\le\sum_{q\ge0}|\cH_q|\,e^{-2\Delta_q^2/N}
\le\sum_{q\ge0}e^{(2q+k+1)(k+1)+q\ln(s+1)-2q(k+2+\ln(s+1))-2(k+2)^2}\\
&=\sum_{q\ge0}e^{-2q-q\ln(s+1)+(k+1)^2-2(k+2)^2}
\le\sum_{q\ge0}e^{-2q-(k^2+6k+7)},
\end{align*}
which is at most $1.2\,e^{-(k^2+6k+7)}\le2^{-(k+6)}$.

For a good table with $U<N$ survivors, the agreement on the $N$ points $u$
is at most $N/2+\alpha\sqrt U+\beta$. This also holds for $U\ge N$, because
then $\alpha\sqrt U\ge N>N/2$. Averaging over all restrictions
and using $\mathbb E\sqrt U\le\sqrt{\mathbb E U}$ (Cauchy--Schwarz) proves
the lemma, since the average agreement over all restrictions equals the
agreement over all inputs.
\end{proof}
Kane and Williams~\cite{KW}, and similarly the later
work~\cite{KKL,CSS}, instead bound $U$
by Markov's inequality and then use a single survivor count. That loses a
factor of $\varepsilon$; the uniform bound with Cauchy--Schwarz does not.

\begin{proof}[Proof of Theorem~\ref{thm:lower}]
We prove more generally, for all $k,m\ge1$ and
$\varepsilon\ge3(k+2)/\sqrt N$, that
\begin{equation}\label{eq:lower-general}
m\,N\,\varepsilon^2\le 36\,\bigl(k+2+\ln(s+1)\bigr)\,s\sqrt n
\quad\text{and}\quad
m^{3/2}N\varepsilon^2\le 36\,\bigl(k+2+\ln(w+1)\bigr)\,w ,
\end{equation}
and that for $m\le N$ the factors $k+2+\ln(s+1)$ and $k+2+\ln(w+1)$ can
both be replaced by $4(k+2)$:
\begin{equation}\label{eq:lower-simple}
mN\varepsilon^2\le144(k+2)\,s\sqrt n
\quad\text{and}\quad
m^{3/2}N\varepsilon^2\le144(k+2)\,w .
\end{equation}
In the balanced case $m=N/k\le N$ and $n=2N$, so \eqref{eq:lower-simple}
reads
\[
s\ge\frac{\varepsilon^2N^{3/2}}{144\sqrt2\,k(k+2)}
\quad\text{and}\quad
w\ge\frac{\varepsilon^2N^{5/2}}{144\,k^{3/2}(k+2)},
\]
which is the theorem, since $N=n/2$ and $k=\log_2(n/2)$.

By hypothesis, $\beta/N=(k+2)/\sqrt N\le\varepsilon/3$ and
$2^{-(k+6)}\le\varepsilon/192$. So Lemma~\ref{lem:avg} gives
$\varepsilon/3\le\alpha\sqrt{\mathbb E[U]}/N$. Squaring and inserting
Lemma~\ref{lem:surv},
\[
\varepsilon^2N^2\le9\alpha^2\,\mathbb E[U]\le
9N\bigl(k+2+\ln(s+1)\bigr)\cdot\frac{4s\sqrt n}{m},
\]
which is the first bound of~\eqref{eq:lower-general}.

For wires, count as survivors only the gates that depend on at least two
live bits. A gate depending on at most one live bit computes that bit, its
negation, or a constant, so it can be absorbed into the top gate, which reads
the live bits anyway. Every remaining survivor has at least two input wires,
so there are at most $w$ candidate gates, and Lemmas~\ref{lem:count}
and~\ref{lem:avg} hold with $s$ replaced by $w$. Now use the second bound of
Lemma~\ref{lem:surv}.

For~\eqref{eq:lower-simple}, let $m\le N$. If $\ln(s+1)\le3(k+2)$, then $k+2+\ln(s+1)\le4(k+2)$.
Otherwise $s\ge e^{3(k+2)}-1>N^{5/2}\ge mN$, so the first bound holds
trivially. The same argument works for $w$, since $N^{5/2}\ge m^{3/2}N$.
\end{proof}

\begin{remark}[Limits of the method]
Theorem~\ref{thm:lower} is within one logarithm of what survivor counting
can give. After a restriction with one live bit per block, the function
$u\mapsto x_{u\oplus o}$ has a depth-two circuit with $N$ bottom gates, and
gates such as $\ind{t_S\ge\theta}$ survive with probability
$\Theta(k/\sqrt N)$; Kane and Williams already observe that their
restriction lemma is tight for majority gates~\cite[Section~3]{KW}. So
counting the survivors of a single restriction cannot give more than about
$n^{3/2}/\log n$ gates.
\end{remark}

%===================================================================
% Copyright (c) 2026 Christian Reitwiessner.
% Authors: Christian Reitwiessner
\section{A near-quadratic lower bound for sign-coherent circuits}\label{sec:coherent}

The upper bound circuit of Theorem~\ref{thm:exact} uses only bottom gates of a special
kind. This section shows that with such gates it cannot be improved by more
than a polylogarithmic factor.

\paragraph*{Sign-coherence.}
A bottom gate is \emph{sign-coherent} if, inside every address block, its
weights are either all $\ge0$ or all $\le0$; its table weights are
arbitrary. A depth-two threshold circuit is \emph{sign-coherent} if all its
bottom gates are; its top gate is arbitrary. Every gate whose address
weights are constant within each block is sign-coherent, as is every gate
with nonnegative address weights. All bottom gates of the circuit in Theorem~\ref{thm:exact} are of this kind.

\begin{theorem}\label{thm:coherent}
Let $C$ be a sign-coherent depth-two threshold circuit for $A$ with $G$
bottom gates. Put $D_0=\max\{0,N-k-1-\log_2(G+1)\}/\bigl(k+\log_2(G+1)\bigr)$. Then
\[
G\ \ge\ \frac{m^k}{m^k-(m-1)^k}\,D_0\ \ge\ \frac mk\,D_0 .
\]
Consequently $G\ge\min\{N^2,\ m(N-3k-1)/(3k^2)\}$, and in the balanced case
\[
G\ \ge\ \frac{N(N-3k-1)}{3k^3}\ =\ \Omega\!\left(\frac{n^2}{\log^3 n}\right).
\]
\end{theorem}

So among sign-coherent circuits the circuit of Theorem~\ref{thm:exact} is
optimal up to a polylogarithmic factor. The earlier circuit of Kane and
Williams~\cite[Thm.~1.2]{KW} is sign-coherent as well. A polynomial-exponent improvement below quadratic therefore requires
gates whose weights have mixed signs inside a block.

\paragraph*{Plan.}
We walk through all unit cubes of the grid of block counts. On each cube,
the circuit reveals the whole table, so many gates must be non-constant there
(Lemma~\ref{lem:entropy}). A sign-coherent gate is non-constant on at most a
$k/m$ fraction of the cubes (Lemma~\ref{lem:diag}). Comparing the two counts
gives the theorem.

\paragraph*{The cube restriction.}
Fix an order of the positions inside every block. For a vector
$s\in\{0,\dots,m\}^k$ let $a(s)$ be the address part in which block $i$ has
ones exactly at its first $s_i$ positions. A \emph{cube} is a vector
$q\in\{0,\dots,m-1\}^k$; there are $m^k$ cubes. For $z\in\bits^k$ the input
$(x,a(q+z))$ differs from $(x,a(q))$ exactly in the live positions
$q_i+1$ of the blocks $i$ with $z_i=1$, and its address is
$(q\bmod2)\oplus z$. Hence, for every cube $q$,
\[
z\mapsto A\bigl(x,a(q+z)\bigr)=x_{(q\bmod 2)\oplus z}
\]
determines the whole table $x$. A bottom gate \emph{splits} the cube $q$ for
the table $x$ if it is non-constant in $z$ on the inputs $(x,a(q+z))$. Let
$d_q(x)$ be the number of bottom gates that split $q$ for $x$.

\begin{lemma}[Many gates split each cube]\label{lem:entropy}
Let $C$ be any depth-two threshold circuit for $A$ with $G$ bottom gates, and
let $x$ be uniformly random. Then $\mathbb E_x[d_q(x)]\ge D_0$ for every cube
$q$.
\end{lemma}
\begin{proof}
Fix $q$ and $x$, and consider bottom gate $j$ as a function of $z$ alone.
It computes $\ind{\lambda_j\cdot z\ge\tau_j(x)}$, where
$\lambda_j\in\mathbb R^k$ lists its fixed weights on the live positions,
and $\tau_j(x)$ absorbs the contributions of the table bits and fixed
address bits.
Let $I$ be the set of gates that split $q$ for this $x$, and write $g_j(z)$
for their restricted functions. Absorb the constant gates, the table bits
and the fixed address bits into the top threshold $\kappa(x)$. The top gate
then computes
$\ind{\sum_{j\in I}v_jg_j(z)+\mu\cdot z\ge\kappa(x)}$, where $v_j$ are its
fixed weights on the bottom gates and $\mu$ lists its fixed weights on the
live positions. So among the weights and thresholds, only $\tau_j$ and
$\kappa$ depend on $x$.

Therefore the function $z\mapsto x_{(q\bmod2)\oplus z}$, and with it $x$, is
determined by: the set $I$ of splitting gates; for each $j\in I$ the
non-constant function $\ind{\lambda_j\cdot z\ge\tau_j(x)}$, one of at most
$2^k-1$ (only the position of $\tau_j(x)$ among the $\le2^k$ values
$\lambda_j\cdot z$ matters, and the two constant functions are excluded);
and the position of $\kappa(x)$. After the restricted bottom-gate
functions have been specified, the top gate's score is fixed at each of the
$2^k$ inputs $z$, so at most $2^k+1$ positions of $\kappa(x)$ matter.
Now let $x$ be uniformly random, and write $H$ for Shannon entropy in bits.
Encoding $|I|=d_q(x)$ first, the entropy chain rule gives
\[
\begin{aligned}
N=H(x)&\le H(d_q(x))+\mathbb E_x\!\left[\log_2\!\left(
\binom{G}{d_q(x)}(2^k-1)^{d_q(x)}(2^k+1)\right)\right]\\
&\le\log_2(G+1)+\mathbb E_x[d_q(x)]\,(\log_2(G+1)+k)+k+1,
\end{aligned}
\]
which rearranges to the claim; a negative numerator is replaced by the
trivial lower bound zero.
\end{proof}

\begin{lemma}[A coherent gate splits few cubes]\label{lem:diag}
For every table $x$, a sign-coherent gate splits at most $m^k-(m-1)^k$ cubes.
\end{lemma}
\begin{proof}
Fix the table $x$ and a sign-coherent gate with threshold $\theta$.
Write $w_{i\ell}$ for its weight on position $\ell$ of block $i$.
On the input $(x,a(s))$, its score is
\[
F(s)=c(x)+\sum_{i=1}^k\sum_{\ell=1}^{s_i}w_{i\ell},
\]
where $c(x)$ is the contribution of the table bits. Let $\varepsilon_i=+1$
if all weights in block $i$ are nonnegative, and $\varepsilon_i=-1$
otherwise. By sign-coherence, moving $s_i$ one step in direction
$\varepsilon_i$ never decreases $F$.

For a cube $q$, define the two opposite vertices
\[
q^-=q+(\mathbf1-\varepsilon)/2,
\qquad q^+=q+(\mathbf1+\varepsilon)/2.
\]
These minimize and maximize $F$ on the cube, respectively. Thus the gate
splits $q$ if and only if $F(q^-)<\theta\le F(q^+)$.

Let $Q=\{0,\dots,m-1\}^k$ be the set of cubes. Partition $Q$ into maximal
\emph{lines} $q,\,q+\varepsilon,\,q+2\varepsilon,\dots$ contained in $Q$.
For consecutive cubes, the identity
\[
(q+\varepsilon)^-=q^+
\]
shows that the largest score on one cube equals the smallest score on the
next. Hence their half-open intervals $(F(q^-),F(q^+)]$ are disjoint
along each line, and the fixed threshold $\theta$ can split at most one
cube per line.

Each line has exactly one starting cube: a $q\in Q$ with
$q-\varepsilon\notin Q$. To count these, note that $q-\varepsilon\in Q$
requires $q_i\in\{1,\dots,m-1\}$ when $\varepsilon_i=+1$, and
$q_i\in\{0,\dots,m-2\}$ when $\varepsilon_i=-1$. There are $m-1$ choices
in each coordinate, hence $(m-1)^k$ cubes with a predecessor.
The remaining $m^k-(m-1)^k$ cubes start lines, proving the bound.
\end{proof}

\begin{proof}[Proof of Theorem~\ref{thm:coherent}]
Sum Lemma~\ref{lem:entropy} over all cubes and compare with
Lemma~\ref{lem:diag}:
$m^kD_0\le\sum_q\mathbb E_x[d_q(x)]\le G\,(m^k-(m-1)^k)$.
Then use $1-(1-1/m)^k\le k/m$. If $G\ge N^2$ there is nothing to prove. Otherwise $\log_2(G+1)\le2k$ and
$D_0\ge(N-3k-1)/(3k)$. In the balanced case $m=N/k$.
\end{proof}

\begin{remark}
(1) Lemma~\ref{lem:diag} is tight for all $k,m$. The fact that a
unate function changes at most once along a chain compatible with its
orientation is classical~\cite[Lemma~2.1]{IPS}; to the best of our knowledge,
the cube restriction and the resulting near-quadratic lower bound are new.
(2) The explicit balanced bound in Theorem~\ref{thm:coherent} first exceeds
$n^{3/2}$ at $k=64$ among the admissible powers of two. The method
(Lemma~\ref{lem:entropy} plus a uniform bound as in Lemma~\ref{lem:diag})
cannot give more than $O(n^2/\log n)$. (3) Lemma~\ref{lem:entropy} holds for
all circuits. For general gates, the limit of survivor counting noted after
Theorem~\ref{thm:lower} does not depend on the choice of restrictions: under
every distribution of restrictions that leaves one live bit per block, a
gate with random $\pm1$ weights is non-constant with probability at least
$1/(2\sqrt{2km})$, so such arguments cannot exceed $O(n^{3/2})$. Gates with mixed signs inside a block are exactly where
the gap between $n^{3/2}$ and $n^2$ lies.
\end{remark}


%===================================================================
\section{Open questions}\label{sec:open}

\begin{enumerate}
\item \emph{Mixed signs.} Is there an exact circuit with $n^{2-\Omega(1)}$
  gates? By Theorem~\ref{thm:coherent} it must use bottom gates whose weights
  have mixed signs inside a block. Conversely, can the exact lower bound beat
  $n^{3/2}$? Such a proof must handle gates whose address weights nearly
  cancel, since those defeat both survival counting and the cube restriction.
  Rectangle methods (discrepancy, sign rank) do not reach $n^{2-o(1)}$
  either: for every split of the inputs between two players, Andreev's
  function has a protocol with $O(k)$ bits of communication, and sign rank
  is too weak against $\THR\circ\THR$ in general~\cite{CM}.
\item \emph{Parity.} Fixing the table to $x_c=y\oplus c_1\oplus\dots\oplus c_k$
  turns a circuit for $A$ into one for parity of $km+1$ bits. With
  sign-coherent gates parity needs about half as many gates as bits; with
  general gates a linear depth-two lower bound for parity is open and tied to
  slicing all edges of the hypercube with hyperplanes. Does a direct-sum
  version, about $N$ times the parity cost, hold for Andreev's function?
\item \emph{Wires.} The exact wire bounds are $\Omega(n^{5/2}/\log^{5/2}n)$
  and $O(n^3)$. Can either exponent move?
\item \emph{From approximation to exactness.} The gate complexity is about
  $n^{3/2}$ for constant error and between $n^{3/2}$ and $n^2$ for error $0$.
  How does it depend on the error $\delta$ between $2^{-\Theta(n)}$ and a
  constant?
\end{enumerate}

\paragraph*{Formalization.}
The Lean~4 proofs are in the \texttt{cslib-frontier} repository at
\url{https://github.com/cslib-community/cslib-frontier}.
We refer to commit\\
\href{https://github.com/cslib-community/cslib-frontier/commit/89ed4de0e54b4953fcc15b8c5d33fad4a344ea84}{\texttt{89ed4de0e54b4953fcc15b8c5d33fad4a344ea84}}\\
on the \texttt{main} branch. The files are in
\href{https://github.com/cslib-community/cslib-frontier/tree/89ed4de0e54b4953fcc15b8c5d33fad4a344ea84/CslibFrontier/Computability/Circuit/Boolean/WeightedThreshold}{\texttt{WeightedThreshold/}}
under \nolinkurl{CslibFrontier/Computability/Circuit/Boolean/}.
The main files are:
\begin{itemize}
\item \texttt{Andreev.lean} and \texttt{AndreevSharp.lean} for the exact
      constructions (Proposition~\ref{prop:basic} and Theorem~\ref{thm:exact});
\item \texttt{AndreevOddErrors.lean} and \texttt{AndreevTruncatedCircuit.lean}
      for Lemma~\ref{lem:odd} and Theorem~\ref{thm:avgupper};
\item \texttt{AndreevLowerBound.lean} for Theorem~\ref{thm:lower};
\item \texttt{AndreevSignCoherent.lean} for Theorem~\ref{thm:coherent} and
      Lemmas~\ref{lem:entropy}--\ref{lem:diag}.
\end{itemize}
All numbered results are covered by these files and their dependencies.
The formalization proves explicit bounds in the parameters; the asymptotic
consequences are derived in the text. The remarks are not formalized.

\paragraph*{Acknowledgements.}
The results, the formalization and this note were developed with AI
assistance.

\appendix
\section{Proof of Theorem~\ref{thm:exact}}\label{app:exact}


\begin{proof}
\emph{Forms used only in this proof.}
Write $s_i=|a_i|$, $\sigma=\sum_i s_i$, and let $e_i$ be the address with
only coordinate $i$ equal to one. Put
\[
z\coloneqq\sum_{c\ne0}x_c,\quad E\coloneqq\sum_i x_{e_i},\quad
w\coloneqq2z+2\alpha E+\sigma,\quad R\coloneqq2(N-1)+2\alpha k+km,\quad C\coloneqq N(R+4).
\]
Thus $0\le w\le R$. For each nonempty $S\subseteq[k]$, define
\[
Y_S\coloneqq\sum_{c:S\cdot c\text{ odd}}x_c,\quad
W_S\coloneqq Y_S+\alpha\sum_{i\in S}x_{e_i},\quad
\kappa_S\coloneqq-4(2\alpha+1)(k-2|S|),
\]
and let $r_S(u)=-8\lceil u/2\rceil^2+\kappa_S\lceil u/2\rceil$.
All of $W_S,t_S,w,s_i$ are affine forms in the original inputs.

\emph{The score.} Consider
\begin{align}\label{eq:bestscore}
P\coloneqq{}&2N+\sum_{S\ne\varnothing}r_S(2W_S+t_S)
 +\frac N2\left(w^2+\sum_i s_i^2\right)\notag\\
 &+4Nz-2N(\alpha^2+1)E-2\alpha N\sigma
 +2N(1+\alpha)\sum_i x_{e_i}s_i-Cx_0\ind{b=0}.
\end{align}
We first verify its sign, then realize its non-affine terms with bottom gates.
The following identities follow by expanding the squares and summing
characters; all sums over $S$ are over nonempty sets:
\begin{align*}
\sum_S4W_S(-1)^{t_S}
 &=2N\ind{b=0}z-2N\ind{b\ne0}x_b
   +2\alpha N\sum_i x_{e_i}\bigl(\ind{b=0}-\ind{b=e_i}\bigr),\\
8\sum_S(W_S^2+W_St_S)
 &=\frac N2(w^2-\sigma^2)+2Nz+(4\alpha+2\alpha^2)NE
   +2N(1+\alpha)\sum_i x_{e_i}s_i,\\
8\sum_S\lceil t_S/2\rceil^2
 &=\frac N2\left(\sigma^2+\sum_i s_i^2\right)
   +N(\sigma+1)\ind{b\ne0}+N\sum_i s_i\ind{b=e_i}.
\end{align*}
For example, the first is character orthogonality applied to $Y_S$ and to
$\ind{i\in S}$; the other two use the same identity and
$\lceil t/2\rceil=(t+\ind{t\text{ odd}})/2$.
Also
\[
4W(-1)^t=4W-8\lceil(2W+t)/2\rceil^2+8W^2+8Wt+8\lceil t/2\rceil^2.
\]
The coefficient shifts $\kappa_S$ supply the address correction because
\[
\sum_S\kappa_S\bigl(\lceil t_S/2\rceil-t_S/2\bigr)
 =N(2\alpha+1)\sum_i\ind{b=e_i}.
\]
Substitution in \eqref{eq:bestscore} now gives the simple closed form
\begin{equation}\label{eq:bestclosed}
P=\begin{cases}
N(w+2)-Cx_0,&b=0,\\
2N(1+\alpha)(1-x_{e_i})-Ns_i,&b=e_i,\\
N(1-2x_b),&\text{otherwise}.
\end{cases}
\end{equation}
At $b=0$, use $0\le w\le R$ and $C=N(R+4)$. At $b=e_i$, the odd
integer $s_i$ lies between $1$ and $2\alpha+1$. In every case $P\ge N$
when $x_b=0$, and $P\le-N$ when $x_b=1$. The top gate outputs $\ind{-P\ge0}$.

\emph{Raw gates and their count.}
The staircase for $r_S(2W_S+t_S)$ uses only odd thresholds
\[
\ind{2W_S+t_S\ge\theta},\qquad
1\le\theta\le N+2\alpha|S|+|S|m,
\]
with coefficient $-8\theta+\kappa_S$. The two square terms use $R+km$
gates by Lemma~\ref{lem:ladder}. For $B=n+2$, each product $x_{e_i}s_i$
uses $m$ gates $\ind{s_i+Bx_{e_i}\ge B+\ell}$, $1\le\ell\le m$.
Finally, the identity
\[
x_0\ind{b=0}=x_0-x_0\sigma+
 \frac4N\sum_{S\ne\varnothing}x_0\lfloor t_S/2\rfloor
\]
implements the last correction with $km$ gates
$\ind{\sigma+Bx_0\ge B+\ell}$ and
$\sum_S\lfloor |S|m/2\rfloor$ gates
$\ind{t_S+Bx_0\ge B+2j}$. Its identity follows from
$\sum_S\ind{t_S\text{ odd}}=(N/2)\ind{b\ne0}$.
The even-level $\sigma$ gates coincide with the $S=[k]$ gates, saving
$\lfloor km/2\rfloor$ gates. Everything else in $P$ is an affine top term.

Since $\sum_{S\ne\varnothing}|S|=kN/2$, the gate counts for the $r_S$
staircases and the last correction sum to
\[
\sum_{S\ne\varnothing}\left(
 \left\lceil\frac{N+2\alpha|S|+|S|m}{2}\right\rceil
 +\left\lfloor\frac{|S|m}{2}\right\rfloor\right)
 =\frac{N(N-1)}2+\frac{\alpha kN}2+\frac{kmN}2.
\]
Adding the remaining gates gives \eqref{eq:bestcount}. Vanishing coefficients
or further coincidences can only improve it. In the balanced case $m$ is even,
so $\alpha=m/2-1$; substitution gives the stated asymptotic bound.
Bottom coefficients and thresholds are $O(n)$; the displayed
expansion has top coefficients bounded by $O(n^3)$.
\end{proof}

\section{Proof of Lemma~\ref{lem:surv}}\label{app:surv}

\begin{lemma}[Erd\H{o}s~\cite{Erdos}]\label{fact:lo}
Let $R\ge1$ be an integer, let $z_1,\dots,z_R$ be independent uniform
bits, and let real weights $\omega_1,\dots,\omega_R$ satisfy $|\omega_j|\ge\mu>0$, and let $I$ be a half-open interval of length
$\mu$. Then
$\Pr\bigl[\sum_j\omega_jz_j\in I\bigr]\le\binom{R}{\lfloor R/2\rfloor}2^{-R}\le1/\sqrt R$.
\end{lemma}

Fix one bottom gate $\ind{\sum_j\omega_jz_j\ge\tau}$, where $j$ runs over
all input positions. Let $\Lambda$ be the set of address positions with
$\omega_j\ne0$, so $|\Lambda|$ is at most the number of the gate's input wires.
For $j\in \Lambda$ let $\rho(j)$ be the number of $j'\in \Lambda$ with
$|\omega_{j'}|\ge|\omega_j|$, including $j$ itself. Fix the table bits; they
only shift the threshold.

\begin{lemma}\label{lem:onerestr}
Fix the live positions. Over the random fixed address bits, the gate is
non-constant on the live bits with probability at most
$2\sum_{j}\rho(j)^{-1/2}$, the sum over the live positions $j\in \Lambda$.
\end{lemma}
\begin{proof}
Let $\Phi$ be the contribution of the fixed bits. The restricted gate is
$\ind{\sum_{j\text{ live}}\omega_ju_j\ge\tau-\Phi}$, and it is non-constant
if and only if $\tau-\Phi$ lies in a half-open interval of length
$\sum_{j\text{ live}}|\omega_j|$. Split this interval into consecutive
half-open pieces of lengths $|\omega_j|$, one for each live $j\in \Lambda$.

For a live $j\in \Lambda$, let $R_j$ be the number of \emph{fixed} positions
$j'\in \Lambda$ with $|\omega_{j'}|\ge|\omega_j|$. If $R_j\ge\rho(j)/4$ for every live
$j\in \Lambda$, then Lemma~\ref{fact:lo}, applied to these $R_j$ positions and
conditioned on all other fixed bits, bounds the probability of piece $j$ by
$R_j^{-1/2}\le2\rho(j)^{-1/2}$, and summing over the pieces proves the lemma.

Otherwise some live $j$ has $R_j<\rho(j)/4$, so more than $3\rho(j)/4$ live
positions $j'\in \Lambda$ have $|\omega_{j'}|\ge|\omega_j|$, and each of them has
$\rho(j')\le\rho(j)$. Then the claimed bound is at least
$\tfrac34\rho(j)\cdot2\rho(j)^{-1/2}\ge\tfrac32$, which exceeds any
probability.
\end{proof}

\begin{proof}[Proof of Lemma~\ref{lem:surv}]
Each position of a block is live with probability $1/m$. Sorting the values
$\rho(j)$, the $t$-th smallest is at least $t$. So the gate survives with
probability at most
\[
\frac2m\sum_{j\in \Lambda}\rho(j)^{-1/2}\le\frac2m\sum_{t=1}^{|\Lambda|}t^{-1/2}
\le\frac{4\sqrt{|\Lambda|}}{m}\le\frac{4\sqrt n}{m}.
\]
Summing over the at most $s$ bottom gates gives the first bound.

A gate that depends on at least two live bits is non-constant and has at
least two live positions in $\Lambda$. Positions in different blocks are live
independently, so for $j\in \Lambda$,
\[
\Pr[\,j\text{ and another position of }\Lambda\text{ are live}\,]
\le\min\Bigl\{\frac1m,\frac{|\Lambda|}{m^2}\Bigr\}\le\frac{\sqrt{|\Lambda|}}{m^{3/2}}.
\]
By Lemma~\ref{lem:onerestr}, the gate depends on at least two live bits with
probability at most
\[
2\sum_{j\in \Lambda}\frac{\sqrt{|\Lambda|}}{m^{3/2}}\,\rho(j)^{-1/2}
\le\frac{2\sqrt{|\Lambda|}}{m^{3/2}}\cdot2\sqrt{|\Lambda|}=\frac{4|\Lambda|}{m^{3/2}}.
\]
Summing over the bottom gates, and using that $\sum|\Lambda|$ is at most the number
of wires $w$, gives the second bound.
\end{proof}

The formal proofs are the Lean theorems
\texttt{card\_not\_forcedToConst\_ltfEval\_le} (Lemma~\ref{lem:onerestr}, there
with all non-live positions random),
\texttt{sum\_circuitUnforcedConstCount\_blockLiveSet\_le} and
\texttt{sum\_circuitUnforcedAtMostOneCount\_blockLiveSet\_le}.

\begin{thebibliography}{HMPST93}

\bibitem[ACW16]{ACW} J.~Alman, T.~M. Chan, and R.~Williams.
\emph{Polynomial Representations of Threshold Functions and Algorithmic
Applications}. FOCS 2016, pages 467--476.
\bibitem[And87]{And87} A.~E. Andreev. \emph{On a method for obtaining more
than quadratic effective lower bounds for the complexity of $\pi$-schemes}.
Moscow Univ.\ Math.\ Bull.\ 42(1):63--66, 1987.
\bibitem[CM17]{CM} A.~Chattopadhyay and N.~S. Mande.
\emph{Weights at the Bottom Matter When the Top is Heavy}.
\href{https://eccc.weizmann.ac.il/report/2017/083/}{ECCC TR17-083}, 2017.
\bibitem[CSS18]{CSS} R.~Chen, R.~Santhanam, and S.~Srinivasan.
\emph{Average-Case Lower Bounds and Satisfiability Algorithms for Small
Threshold Circuits}. Theory of Computing 14(9):1--55, 2018.
\bibitem[CTW26]{CTW} L.~Chen, A.~Tal, and Y.~Wang.
\emph{Superquadratic Lower Bounds for Depth-2 Linear Threshold Circuits}.
STOC 2026, pages 874--885. \href{https://eccc.weizmann.ac.il/report/2026/039/}{ECCC TR26-039}.
\bibitem[Erd45]{Erdos} P.~Erd\H{o}s. \emph{On a lemma of Littlewood and Offord}.
Bull.\ Amer.\ Math.\ Soc.\ 51:898--902, 1945.
\bibitem[HMPST93]{HMPST} A.~Hajnal, W.~Maass, P.~Pudl\'ak, M.~Szegedy, and
G.~Tur\'an. \emph{Threshold circuits of bounded depth}. J.\ Comput.\ Syst.\
Sci.\ 46(2):129--154, 1993.
\bibitem[H{\aa}s98]{Has98} J.~H{\aa}stad. \emph{The shrinkage exponent of
de Morgan formulas is $2$}. SIAM J.\ Comput.\ 27(1):48--64, 1998.
\bibitem[IPS97]{IPS} R.~Impagliazzo, R.~Paturi, and M.~Saks. \emph{Size--depth tradeoffs for threshold circuits}. SIAM J.\ Comput.\ 26(3):693--707, 1997.

\bibitem[KKL17]{KKL} V.~Kabanets, D.~M. Kane, and Z.~Lu.
\emph{A Polynomial Restriction Lemma with Applications}. STOC 2017, pages
615--628.
\href{https://eccc.weizmann.ac.il/report/2017/026/}{ECCC TR17-026}.
\bibitem[KW16]{KW} D.~M. Kane and R.~Williams.
\emph{Super-Linear Gate and Super-Quadratic Wire Lower Bounds for Depth-Two
and Depth-Three Threshold Circuits}. STOC 2016, pages 633--643.
\href{https://arxiv.org/abs/1511.07860}{arXiv:1511.07860v1}.
\bibitem[Nag26]{Nag} D.~Nag. \emph{Almost-Everywhere Near-Cubic Wire Lower Bounds for
$\mathsf{SYM}\circ\THR$ and $\THR\circ\THR$}. Preprint,
\href{https://eccc.weizmann.ac.il/report/2026/167/}{ECCC TR26-167}, 2026.
\bibitem[Nis93]{Nis93} N.~Nisan. \emph{The communication complexity of
threshold gates}. In Combinatorics, Paul Erd\H{o}s is Eighty, Vol.~1, pages
301--315, 1993.
\bibitem[RSO94]{RSO} V.~P. Roychowdhury, K.-Y. Siu, and A.~Orlitsky.
\emph{Neural models and spectral methods}. In Theoretical Advances in Neural
Computation and Learning, pages 3--36, 1994.
\bibitem[Tal14]{Tal14} A.~Tal. \emph{Shrinkage of De Morgan formulae by
spectral techniques}. FOCS 2014, pages 551--560.
\bibitem[Tam16]{Tam} S.~Tamaki. \emph{A Satisfiability Algorithm for Depth Two
Circuits with a Sub-Quadratic Number of Symmetric and Threshold Gates}.
\href{https://eccc.weizmann.ac.il/report/2016/100/}{ECCC TR16-100}, 2016.
\end{thebibliography}
\end{document}
